\section{Introduction}
\label{sec:intro}

Assistants that talk with a user over months need a long-term memory of these conversations. A line of recent studies finds that the simplest memory, which keeps the raw history and retrieves the rounds most similar to a question, matches or beats systems that summarize or restructure the history in advance~\citep{bao2026rawrag,zhou2026explicit,sharma2026selection,peng2026memorytax}. On text benchmarks it is nearly as accurate as reading the whole history: in our measurements, cosine retrieval comes within one point of full-context reading on LoCoMo and LongMemEval-S, where it answers 94\% of the questions (Section~\ref{sec:problem}).

We find that this conclusion does not carry over to multimodal memory, although more and more of what users share is visual, such as receipts, screenshots, product photos, and charts. With a strong reader, the same retrieval falls 7 points below full-context reading on MemEye~\citep{guo2026memeye}, whose answers depend on images, and 16 points below it on VisualMem~\citep{nguyen2026personalvisual}, where it returns the image that holds the answer for fewer than half of the questions. The cause lies in the questions. Many ask about a \emph{set} of past items, such as every advertisement of one brand, the summary cards of two particular days, or each checkpoint of a route, and the answer depends on every member. Similarity retrieval scores each item on its own: the 10 most similar rounds hold the complete evidence for 92\% of single-round questions but for only 32\% of questions with four or five evidence rounds (Figure~\ref{fig:setsize}). Set members are also often named by references to the history, such as ``the D9 summary'' for a card shared in the ninth session, which neither the round nor its caption contains. We call this failure the \emph{set-completeness gap}.

We argue that finding such evidence is a decision about a set, and that this decision is best made by a model that sees the whole history at once. Our approach separates navigation from perception. The memory keeps every round raw and adds one short index line per round. At question time, a model reads the index, much as a reader scans a table of contents, and returns the rounds that form the evidence set, choosing how many it needs. This navigation has a simple precondition: the model can only find evidence that the index shows. We measure this precondition as \emph{index visibility}, the share of the question words in the evidence that the index lines expose, and estimate it at question time without annotations.

Visibility-Gated Set Reading (\vgs{}) turns this precondition into one rule, which we apply unchanged to every benchmark and reader (Figure~\ref{fig:overview}). When the estimated visibility is high, a selector returns the evidence set from the index, and the reader answers from the original rounds and images. When it is low, \vgs{} widens the reading in the way that the reader needs: a lightweight reader, which often fails to combine the members of a set, receives wider index lines and one explicit note per member, whereas a strong reader reads the rounds most similar to the question.

On three multimodal benchmarks, \vgs{} improves a lightweight reader over top-40 retrieval by 5.3 points on MemEye, 7.1 points on MemLens~\citep{ren2026memlens}, and 11.5 points on VisualMem, and a strong reader by 4.8 and 3.2 points on MemEye and VisualMem, where the index exposes the evidence, while matching retrieval on MemLens. Six recent memory systems fall below top-40 retrieval on all three benchmarks, so the advantage of raw retrieval over memory systems holds for images as well, but raw retrieval itself leaves the sets incomplete. Full-context reading recovers much of the gap for the strong reader at up to 18 times the cost of retrieval, but not for the lightweight reader, which \vgs{} helps most. Ablations locate the gain: a model that reads the index finds complete sets that similarity misses, hiding the evidence from the index largely removes this advantage, and the gate chooses the better reading on every benchmark, whereas each fixed reading loses up to seven points on one of them.

Our contributions are fourfold. First, we show that the strength of raw retrieval in text memory does not carry over to multimodal memory and trace the gap to the completeness of evidence sets. Second, we propose \vgs{}, a training-free memory reader that navigates an index of the history and returns evidence sets, with one gate for all benchmarks and readers. Third, we identify index visibility as the condition under which navigation succeeds, with causal tests on two benchmarks. Fourth, we show that navigation stays practical: a cache-native index keeps \vgs{} about as fast as retrieval, and compact index lines nearly halve its cost on long histories, whereas bounding the index by similarity brings the failure back.

\begin{figure*}[t]
\centering
\resizebox{\textwidth}{!}{% Figure 1: overview of Visibility-Gated Set Reading (TikZ, 16 cm wide; included by sections/01_intro.tex)
\begin{tikzpicture}[yscale=1.13,
  font=\fontsize{6.9}{8.2}\selectfont, >={Stealth[length=1.8mm,width=1.5mm]},
  panel/.style={rounded corners=3pt, line width=0.6pt},
  title/.style={font=\fontsize{8.1}{9.5}\selectfont\bfseries, anchor=west, inner sep=0pt},
  num/.style={circle, fill=black!78, text=white, font=\fontsize{5.6}{6}\selectfont\bfseries, inner sep=0pt, minimum size=8pt},
  body/.style={align=left, text width=#1, inner sep=0pt, anchor=north west},
  arr/.style={->, line width=0.8pt, draw=black!70},
  chip/.style={rounded corners=1.2pt, draw=gold!70!black, fill=gold!35, font=\ttfamily\fontsize{5.2}{6}\selectfont, inner sep=1.4pt},
  idx/.style={font=\ttfamily\fontsize{4.6}{5.4}\selectfont, anchor=west, inner sep=0pt, text=black!85},
  note/.style={font=\fontsize{6.0}{7.0}\selectfont\itshape, text=black!62},
]
\colorlet{cMem}{blue!6}\colorlet{cMemD}{blue!50!black}
\colorlet{cGate}{yellow!15}\colorlet{cGateD}{orange!60!black}
\colorlet{cStd}{green!8}\colorlet{cStdD}{green!38!black}
\colorlet{cWid}{orange!10}\colorlet{cWidD}{orange!65!black}
\colorlet{cRead}{violet!7}\colorlet{cReadD}{violet!55!black}
% ---------------------------------------------------------------- phases
\draw[black!40, line width=0.45pt] (0.05,6.02) -- (3.9,6.02);
\node[note, fill=white, inner sep=1.5pt] at (1.97,6.02) {write time};
\draw[black!40, line width=0.45pt] (4.35,6.02) -- (15.95,6.02);
\node[note, fill=white, inner sep=1.5pt] at (10.15,6.02) {question time};
\node[note, text=cStdD] at (9.5,5.62) {navigation over the index};
\node[note, text=cReadD] at (14.05,5.62) {perception of the raw rounds};
% ---------------------------------------------------------------- (1) memory with an index
\draw[panel, draw=cMemD, fill=cMem] (0,0) rectangle (3.95,5.4);
\node[num] at (0.25,5.13) {1}; \node[title] at (0.42,5.13) {Memory with an index};
\foreach \y/\t in {4.55/{R37 \textperiodcentered{} 05-06 \textperiodcentered{} dinner at the market}, 4.05/{R12 \textperiodcentered{} 05-03 \textperiodcentered{} booked the hotel}} {
  \draw[rounded corners=1.5pt, draw=black!40, fill=white] (0.22,\y-0.2) rectangle (3.73,\y+0.2);
  \draw[draw=black!40, fill=blue!15] (0.3,\y-0.14) rectangle (0.62,\y+0.14);
  \fill[blue!45] (0.32,\y-0.12) -- (0.43,\y+0.02) -- (0.51,\y-0.05) -- (0.6,\y+0.07) -- (0.6,\y-0.12) -- cycle;
  \fill[yellow!70!orange] (0.55,\y+0.09) circle (0.028);
  \node[idx] at (0.7,\y+0.07) {\t};
  \fill[black!16] (0.7,\y-0.11) rectangle (3.5,\y-0.06);
}
\node[note] at (1.97,3.6) {raw rounds and images, kept unchanged};
\draw[arr, draw=cMemD] (1.97,3.45) -- node[right, note, xshift=1pt] {one line per round} (1.97,3.05);
\draw[rounded corners=2pt, draw=cMemD!75, fill=white, line width=0.5pt] (0.18,0.18) rectangle (3.77,2.95);
\foreach \y/\t in {
  2.70/{R12 05-03 booked the hotel, Alfama},
  2.50/{\hphantom{R12 05-03 }IMG: hotel invoice, 3 nights},
  2.20/{R20 05-04 museums open on Monday?},
  1.90/{R37 05-06 dinner at the market hall},
  1.70/{\hphantom{R37 05-06 }IMG: restaurant bill},
  1.40/{R44 05-07 tram 28 or a walking tour?},
  1.10/{R58 05-09 the D9 summary card},
  0.90/{\hphantom{R58 05-09 }IMG: day summary, receipt},
  0.60/{R61 05-09 taxi to the airport},
  0.40/{\hphantom{R61 05-09 }IMG: taxi receipt}} {
  \node[idx] at (0.42,\y) {\t};
}
\foreach \y in {2.60, 1.80, 1.00, 0.50} { \fill[gold!85!black] (0.3,\y) circle (0.04); }
% ---------------------------------------------------------------- (2) question and visibility gate
\draw[panel, draw=cGateD, fill=cGate] (4.35,0) rectangle (6.95,5.4);
\node[num] at (4.6,5.13) {2}; \node[title] at (4.77,5.13) {Visibility gate};
\node[draw=black!50, fill=white, rounded corners=2.5pt, align=center, text width=2.15cm, inner sep=2.5pt,
      font=\fontsize{6.2}{7.4}\selectfont\itshape] (Q) at (5.65,4.38) {Q: How much did I spend on all receipts from my Lisbon trip?};
\node[draw=cGateD!70, fill=white, rounded corners=2pt, align=left, text width=2.15cm, inner sep=2.5pt] (V) at (5.65,2.85)
      {$\hat v(q)$: share of the question words in the 10 most similar rounds that their index lines show};
\node[draw=cGateD, fill=white, diamond, aspect=1.7, inner sep=0.3pt, line width=0.6pt] (D) at (5.65,1.25) {$\hat v(q)\ge\tau$\,?};
\draw[arr] (Q.south) -- (V.north);
\draw[arr] (V.south) -- (D.north);
% ---------------------------------------------------------------- (3) standard path
\draw[panel, draw=cStdD, fill=cStd] (7.35,2.85) rectangle (11.75,5.4);
\node[num] at (7.6,5.13) {3}; \node[title] at (7.77,5.13) {Standard path: set selection};
\node[body=4.1cm] at (7.5,4.85) {A selector reads the index lines of the \emph{whole} history and returns the evidence set $S(q)$, choosing its size for each question.};
\tikzset{gchip/.style={chip, draw=black!35, fill=black!8, text=black!55}}
\node[anchor=west, font=\fontsize{5.8}{6.8}\selectfont] at (7.5,3.48) {cosine top-4:};
\node[chip] at (9.0,3.48) {R12}; \node[gchip] at (9.55,3.48) {R20}; \node[gchip] at (10.1,3.48) {R44}; \node[chip] at (10.65,3.48) {R37};
\node[note, anchor=west] at (10.95,3.48) {2 of 4};
\node[anchor=west, font=\fontsize{5.8}{6.8}\selectfont\bfseries] at (7.5,3.1) {set selection:};
\node[chip] at (9.0,3.1) {R12}; \node[chip] at (9.55,3.1) {R37}; \node[chip] at (10.1,3.1) {R58}; \node[chip] at (10.65,3.1) {R61};
\node[note, anchor=west, text=cStdD] at (10.95,3.1) {all 4};
% ---------------------------------------------------------------- (4) widened path
\draw[panel, draw=cWidD, fill=cWid] (7.35,0) rectangle (11.75,2.45);
\node[num] at (7.6,2.18) {4}; \node[title] at (7.77,2.18) {Widened path: reader-aware};
\draw[rounded corners=2pt, draw=cWidD!70, fill=white] (7.5,0.98) rectangle (11.6,1.9);
\node[body=3.9cm] at (7.6,1.8) {\textbf{Lightweight reader:} wider index lines (600 characters) and one explicit note per member of the set.};
\draw[rounded corners=2pt, draw=cWidD!70, fill=white] (7.5,0.15) rectangle (11.6,0.83);
\node[body=3.9cm] at (7.6,0.73) {\textbf{Strong reader:} the 40 rounds most similar to the question.};
\draw[arr] (D.east) -- (7.15,1.25) |- (7.35,4.1);
\node[font=\fontsize{5.6}{6}\selectfont, anchor=south] at (6.78,1.27) {yes};
\draw[arr] (D.south) |- (7.35,0.62);
\node[font=\fontsize{5.6}{6}\selectfont, anchor=south] at (6.4,0.63) {no};
% ---------------------------------------------------------------- (5) reader
\draw[panel, draw=cReadD, fill=cRead] (12.15,0) rectangle (16,5.4);
\node[num] at (12.4,5.13) {5}; \node[title] at (12.57,5.13) {Reader};
\node[body=3.6cm] at (12.3,4.85) {Opens the selected rounds and answers from their original text and images.};
\foreach \x/\lab in {12.42/R12, 13.32/R37, 14.22/R58, 15.12/R61} {
  \draw[draw=black!40, fill=white, rounded corners=1pt] (\x,3.32) rectangle (\x+0.72,3.98);
  \fill[blue!15] (\x+0.06,3.38) rectangle (\x+0.66,3.82);
  \fill[blue!45] (\x+0.08,3.4) -- (\x+0.27,3.64) -- (\x+0.38,3.52) -- (\x+0.63,3.78) -- (\x+0.63,3.4) -- cycle;
  \node[font=\ttfamily\fontsize{5.2}{6}\selectfont, text=black!70] at (\x+0.36,3.17) {\lab};
}
\draw[rounded corners=2pt, draw=cReadD, fill=white, line width=0.6pt] (12.32,0.18) rectangle (15.83,2.45);
\node[align=left, text width=3.3cm, inner sep=0pt] at (14.07,1.32) {\textbf{Answer} that combines every member of the set: the hotel invoice, the restaurant bill, the receipt on the D9 summary card, and the taxi receipt.};
\draw[arr] (14.07,2.98) -- (14.07,2.5);
\draw[arr] (11.75,4.1) -- (12.15,4.1);
\draw[arr] (11.75,1.25) -- (12.15,1.25);
\end{tikzpicture}
}
\caption{\textbf{Overview of Visibility-Gated Set Reading (\vgs{}).} At write time, the memory keeps raw rounds and writes one index line per round (gold markers: the evidence set of the example question). At question time, the gate estimates from the question and its most similar rounds whether the index exposes the evidence. On the standard path, a selector reads the index of the whole history and returns the evidence set, whereas cosine retrieval of the same size returns neighbors and misses members; on the widened path, the reading adapts to the reader. The reader always answers from the original rounds and images.}
\label{fig:overview}
\end{figure*}
